% Approach (draft, Week 17; written 4 Oct 2026 from docs/v2/THEORY.md, PROTOCOL.md v1.0 and
% vmsim/policies/{bam,guard}.py). Numbers from results/v2/theory_check.json.
% Only statements that survived the numerical check are here (THEORY.md section 2).

\section{Approach}
\label{sec:approach}

\paragraph{The decision.}
Each day the operator may move VMs off hosts it expects to fail, within a budget $B$ of memory
moved per day. Moving VM $i$ off host $j$ averts an interruption of cost $c_i$ if $j$ fails within
the horizon $H$ and the migration finishes first, and costs $m_i$ either way. With $p_j$ the
probability that $j$ fails within $H$ and $e_i$ the probability that the migration finishes in time,
the expected benefit is $b_i = p_j c_i e_i - m_i$ and the weight is $w_i$, the VM's memory.

\paragraph{Budget-aware migration (BAM).}
BAM ranks VMs by density $b_i / w_i$ and moves them while $b_i > 0$ and the budget lasts. Two
details matter in practice. The interruption cost is weighted by the probability that the VM is
still running $H$ days ahead given its age, since most VMs leave long before their host fails. And
$e_i$ depends on the VM's place in its host's migration queue, so it is recomputed after every
pick. For fixed $e$, choosing what to move is a knapsack, and density order is optimal for its
relaxation and within one item of the integer optimum~\cite{dantzig1957discrete}. On simulator
snapshots the greedy is within 0.05\% of an exact solver on average.

\paragraph{What the theory does and does not promise.}
Under four assumptions (calibrated $p$, failures independent given $p$, the same budget for every
policy, additive costs) and accounting over the horizon the decision used, BAM is never worse in
expectation than a threshold on $p$ given the same budget. That accounting rule is not a detail. v1
credited migrations with failures months later and found the opposite sign.

Under a Gaussian copula on $(\log p, \log\kappa)$ with $\kappa = c e / w$, the advantage has a
closed form, which a Monte Carlo confirms at all 3,024 points of our grid. Its main term is
$\Phi(z_q + \sigma_z) - \Phi(z_q + \sigma_p + \rho\,\sigma_\kappa)$, which is never negative and is
zero only when costs do not vary or are perfectly aligned with risk. The advantage is therefore
driven by the spread of costs, and is largest near $\rho = 0$, not above a break-even correlation.

Two cautions come from the same analysis.
\begin{itemize}
  \item Noise in $p$ does not by itself make BAM lose when costs are not anti-aligned with risk.
  \item Flattening miscalibration ($p^{\gamma}$ with $\gamma < 1$) can make BAM lose, because its
  stopping rule compares a distorted $p\,c$ with $m$.
\end{itemize}

\paragraph{Safe-adaptive guard (P6).}
Because the promise rests on calibration, P6 runs the threshold policy unless two conditions hold
at two consecutive weekly looks.
\begin{enumerate}
  \item \emph{Model-based advantage.} The model-based advantage of BAM's plan over the threshold's
  plan has a positive lower confidence bound. The bound comes from a host bootstrap, with the
  error budget spent as $\alpha_k = 6\alpha/(\pi^2 k^2)$ so that repeated looks stay valid.
  \item \emph{Calibration audit.} A reliability slope fitted on the outcomes that matured in the
  last 90 days lies in $[0.8, 1.25]$.
\end{enumerate}
P6 returns to the threshold policy at the first look that fails either test, and logs every switch
with its evidence. The plug-in advantage uses BAM's own benefit formula, so it is optimistic by
construction. The audit is what protects the operator: on a control whose scores carry no
information, P6 never switched.
